% =====================================================================
%  Section: Discussion and Open Challenges
% =====================================================================
\section{Discussion and Open Challenges}
\label{sec:discussion}

\textbf{Content leakage and style--content disentanglement.} Image-trained stylization
often transfers semantic content from the style reference, distorting the target.
Recent works~\cite{ye2025stylemaster,chen2024artadapter} tackle this through patch
filtering and auxiliary content adapters, but the problem remains open for long videos.

\textbf{Long-video temporal coherence.} Sliding-window strategies~\cite{li2026echostyle}
and content deformation fields~\cite{ouyang2024codef} mitigate drift but struggle with
large motion and complex scenes. Cross-turn consistency in iterative editing is a
related open problem.

\textbf{Art-specific aesthetic preservation.} General style transfer fails to capture
art-specific features such as ink density, brush direction, and compositional blank
space. Specialized losses and datasets~\cite{he2023bitgan,liu2026chipdiff,zhang2025inkspirit}
are essential but currently limited to specific art traditions.

\textbf{Integration of generation and animation pipelines.} Combining painting-specific
style LoRAs, ControlNet conditioning, I2V diffusion models, and talking-portrait
animation into a single pipeline---as exemplified by the OilPaintingVideo project---is
an emerging engineering direction with significant educational potential.

\textbf{Open-source local video generation backbones.} Recent open-weights video
diffusion models such as LTX-Video~2.x and Wan2.1/2.2-I2V provide powerful, locally
deployable backbones for famous-painting video generation. They offer controllable
image-to-video synthesis with manageable GPU footprints, but are typically trained on
natural imagery and struggle to preserve the rich textures and stylistic features of
artworks. Bridging this domain gap---through style-aware fine-tuning, art-specific
conditioning, or stroke-level guidance---remains an open and practically important
problem.

\textbf{Human-in-the-loop artistic guidance.} Existing pipelines accept textual or
image references as guidance, but professional artistic creation typically involves
iterative human input through sketches, color blocks, and line drawings. Allowing
artists to steer video generation with hand-drawn drafts at multiple abstraction
levels---from precise line constraints to loose color palettes---is largely unexplored
and offers a promising direction for human--AI co-creation.
